% Person 3 report section for Milestone 2.

\section{Document Model and Indexing}
\label{sec:indexing}

The final collection contains 6{,}104 movies and 12{,}057 reviews. Each movie is
one JSON Lines document; reviews are embedded in the source but indexed as child
documents. In Solr, every movie is a parent with \texttt{doc\_type=movie} and an id
of the form \texttt{movie-<movie\_id>}, and every review is a child with
\texttt{doc\_type=review} and an id of the form \texttt{movie-<movie\_id>} plus
\texttt{-review-<n>}. The index therefore holds 18{,}161
documents. This parent/child model avoids duplicating movie text for every review
and lets a review query return one movie ranked by its best matching review;
concatenating all reviews into a single field would lose per-review scoring. Only
3{,}544 movies (58.1\%) have at least one review, which limits the reach of
review-based needs.

\subsection{Indexing process}
\label{sec:indexing-process}

\texttt{scripts/build\_collection.py} joins the five interim sources using the
numeric TMDB identifier (\texttt{movie\_id}), rejects non-numeric ids, requires a
Wikipedia plot of at least 50 words, strips footnote markers, and validates every
document against
\texttt{docs/collection.schema.json}. The final run dropped 3{,}231 movies without
a usable plot and 374 with fewer than 50 plot words, producing 6{,}104 documents
and 12{,}057 reviews. The script writes \texttt{collection.jsonl}, a manifest with
counts, filters, build date, and git commit, and \texttt{solr\_documents.json}.
The converter flattens \texttt{cast}, \texttt{crew\_roles}, and \texttt{wikipedia}
into Solr fields and normalizes review dates to full timestamps.
\texttt{solr/startup.sh} recreates a \texttt{solr:10} container, applies
\texttt{solr/schema.json}, uploads the documents with \texttt{commit=true}, and
checks that the indexed parent and child counts match the collection.

\subsection{Fields and analyzer}
\label{sec:fields}

Table~\ref{tab:index-fields} summarizes the indexing roles. Free-text fields use
\texttt{textFoldLower}: StandardTokenizer, ASCIIFolding, LowerCase, English
stopwords, and KStem. The analyzer makes search case- and accent-insensitive
(\textit{Amélie} $\rightarrow$ \textit{amelie}), removes stopwords
(\textit{the}, \textit{and}), and stems some variants
(\textit{investigating} $\rightarrow$ \textit{investigate}). Language and status
are exact \texttt{string} fields; year, runtime, ratings, budget, and revenue are
numeric; release and review dates use \texttt{pdate}. \texttt{crew\_roles.job} is
deliberately display-only: indexing job words would make ``a movie about a
producer'' return films with the most producer credits. Exact docValues queries
remain possible (\texttt{crew\_roles.job:Producer}), but the field contributes no
full-text relevance. \texttt{cast.character} is searchable because users remember
characters, not only actors.


\begin{table}[t]
\centering
\setlength{\tabcolsep}{4pt}
\caption{Indexing roles in the final Solr schema.}
\label{tab:index-fields}
\footnotesize
\begin{tabular}{p{1.9cm}p{5.0cm}}
\hline
Role & Fields \\
\hline
Full text & \texttt{title}, \texttt{overview}, \texttt{tagline}, \texttt{plot},
\texttt{genres}, \texttt{director}, \texttt{cast.name},
\texttt{cast.character}, \texttt{crew\_names}, \texttt{wikipedia.title},
\texttt{review\_text} \\
Exact/numeric & \texttt{year}, \texttt{runtime},
\texttt{original\_language}, \texttt{status}, \texttt{vote\_average},
\texttt{vote\_count}, \texttt{popularity}, \texttt{budget}, \texttt{revenue},
\texttt{release\_date}, \texttt{num\_reviews} \\
Display only & \texttt{crew\_roles.name}, \texttt{crew\_roles.job},
\texttt{cast.order}, \texttt{wikipedia.url}, \texttt{wikipedia.wikidata\_id},
\texttt{wikipedia.imdb\_id}, \texttt{review\_author}, \texttt{review\_rating},
\texttt{review\_created\_at} \\
\hline
\end{tabular}
\end{table}

\subsection{Verification and demo}
\label{sec:verification-demo}

\texttt{scripts/verify\_index.py} checks the manifest and collection counts, a
sample parent and its review children, a parent without children, the analyzer,
sanity queries, and the display-only fields. It reports 25 passed checks and two
documented docValues warnings. The live demo starts the index with
\texttt{bash solr/startup.sh}, shows the counts, the Schema and Analysis screens,
runs the verification, and compares the two retrieval setups on one plot need and
one review need with \texttt{scripts/demo\_queries.py}. The full walkthrough, expected
outputs, and troubleshooting are in \texttt{docs/demo.md}. The pipeline is
reproducible from a clean checkout: \texttt{build\_collection.py} rebuilds the
collection and Solr upload file, \texttt{startup.sh} recreates the index and checks
the counts, and the manifest records the build date and git commit.

\paragraph{Limitations.} Review-based needs can only reach the 3{,}544 movies with
at least one review (58.1\%), so their recall is bounded by coverage rather than
query quality. Plot-based needs depend on Wikipedia coverage and on the 50-word
filter, which excluded 3{,}605 of 9{,}709 movies. The index is lexical: a query
must share vocabulary with the plot or review text, which motivates the semantic
retrieval planned for Milestone 3.
